% ===== §30 =====
\section{Constrained but Not Determined}
\label{p24:sec:constraint}

\noindent
The previous section reached the darkest point of the argument, the completion of alienation in the closed loop of capture. This section is the turn. It does not undo the diagnosis or soften it; it locates the diagnosis correctly, and the correct location is itself the ground of a recovery. The turn rests on the single dialectical principle the series has maintained throughout, and which the first part established for culture in general: that material conditions constrain but do not determine. Applied to the capture just described, this principle yields two consequences that together open the way from diagnosis to ethics, and it is the office of this brief but pivotal section to draw them.

\subsection*{The capture is a historical form, and its alienation is material}

The first consequence is that the capture of intimate culture, however deep, is a \emph{historical form} and not an eternal condition. It is the specific shape that the production and capture of intimate culture take under the present relations of production, and its alienation is therefore material and historical---produced by a particular arrangement of ownership and medium, and, like every such arrangement, particular to its conditions rather than written into the nature of intimate relation as such. This is not a consolation offered against the diagnosis but a precise dating of it. The alienation the previous section described is real, and it is the real form of intimate culture under present conditions; but to see it as a historical form is to see that it is not the truth of intimate culture for all time, not a fate inscribed in the structure of the bond, but the effect of a determinate material arrangement that had a beginning and can have an end. The disciplined historical eye the fourth revision required is turned here upon the present itself: the capture is to be evaluated as a form produced under its conditions, which means it is neither to be mistaken for the permanent nature of intimacy nor to be treated as a doom against which nothing can be done, but understood as what it is, a historical form of alienation, material in its origin and therefore material in its possible supersession.

\subsection*{The capacity for endogenous generation is never wholly destroyed}

The second consequence is the decisive one, and it follows from the framework's account of the Real. The capture colonizes the couple's generation, supplies its forms, at the limit delegates its generative act; but it cannot wholly destroy the capacity for endogenous generation, because that capacity is rooted in something the capture cannot reach. The second revision secured that the relational Real is not completely symbolizable, and the capture operates by symbolic means---by supplying symbolic forms, scripts, and now generative processes, all of which are of the order of the symbolic. What lies in the register of the Real---the unsymbolizable core of the bond, the ``us'' that no supplied script and no generative machine can compose, because it is not a symbolic content to be composed but the real coherence of a particular coupling---remains beyond the reach of a capture that works by symbolic supply. However completely the forms of a couple's intimate culture are colonized, the relational Real at the core of the bond is not colonized, because it is not of the order the colonization operates on; and so long as it remains, the couple retains the ground from which an endogenous generation could begin again, the anchor in the Real from which a new symbolic form, genuinely their own, could be generated in place of the supplied one. The capacity for endogenous generation is never wholly destroyed because its root is in a register the capture cannot symbolize and therefore cannot own.

\subsection*{The negation of the negation}

Read through the Hegelian line, these two consequences are the second negation, and naming them as such completes the dialectical figure the paper has been building. Alienation was the first negation: the exteriorized culture, going out of the couple into a made world, was captured and turned against them, standing over against its makers as an alien power---the negation of the couple's generativity by its own product, seized by another. The turn of this section is the negation of that negation. It is not a return to an imagined state before capture, an innocence of intimate culture that never existed and could not be recovered; it is a determinate negation of the alienation, which carries the couple forward, not back, to a generation that has passed through capture and reclaimed the generative act from it. The exteriorized culture that had become an alien power can be reclaimed, taken up again by the relation whose product it was, and returned to the movement of the couple's self-knowledge from which the capture had diverted it. This is why the diagnosis of alienation is not the paper's last word, and cannot be: because alienation, in the dialectical account, is a moment and not a terminus, the first negation that calls forth its own negation, and the reclamation of the captured culture is the second negation that the framework not only permits but requires. The material conditions constrain, and under present conditions the constraint takes the form of a deep capture; but the conditions do not determine, and the undetermined remainder---the relational Real the capture cannot reach, the endogenous capacity it cannot destroy---is the ground on which the negation of the negation stands. What that reclamation asks of the couple, in the way of a cultural practice and an ethics, is the subject of the section that follows, which turns from the conditions of recovery to the practice of it.
